% BHU PYQ -- Physical Education: Health, Fitness & Exercise Management (2024-25 NEP)
\documentclass[11pt,a4paper]{article}
\usepackage[a4paper,top=2.5cm,bottom=2.5cm,left=2.8cm,right=2.8cm,headheight=14pt]{geometry}
\usepackage{amsmath,amssymb}
\usepackage[T1]{fontenc}
\usepackage[utf8]{inputenc}
\usepackage{lmodern,microtype,fancyhdr,enumitem,hyperref,lastpage,parskip}

\hypersetup{colorlinks=true,urlcolor=black,linkcolor=black,pdftitle={PEDMD11: Health, Fitness & Exercise Management -- BHU NEP 2024-25}}

\pagestyle{fancy}\fancyhf{}
\renewcommand{\headrulewidth}{0.4pt}\renewcommand{\footrulewidth}{0.4pt}
\fancyhead[L]{\small \textit{Banaras Hindu University}}
\fancyhead[R]{\small \textit{PEDMD11: Health \& Fitness (2024-25)}}
\fancyfoot[C]{\small Page \thepage\ of \pageref{LastPage}}

\newlist{parts}{enumerate}{2}
\setlist[parts,1]{label=(\alph*),leftmargin=*,itemsep=4pt,topsep=3pt}
\setlist[parts,2]{label=(\roman*),leftmargin=*,itemsep=2pt,topsep=2pt}

\newcommand{\pts}[1]{\hfill \textit{[#1]}}

\begin{document}

\begin{center}
  {\large\bfseries BANARAS HINDU UNIVERSITY}\\[2pt]
  {\normalsize Under Graduate Multidisciplinary Course Examination, 2024-25 (NEP)}\\[6pt]
  \rule{0.8\linewidth}{0.5pt}\\[6pt]
  {\large\bfseries PHYSICAL EDUCATION \& HEALTH}\\[4pt]
  {\large\bfseries Paper Code: PEDMD11 : Health, Fitness \& Exercise Management}\\[4pt]
  {\normalsize Time: 2:30 Hours \hfill Full Marks: 70}\\[2pt]
  {\small REF NO. CE/Conf./D-I/112}
\end{center}
\rule{\linewidth}{0.4pt}
\smallskip

\noindent \textit{Instructions: Attempt all Sections A, B and C according to directions. Write your Roll No.\ at the top immediately on receipt of this question paper.}

\bigskip

\section*{SECTION -- A}
\noindent \textit{Note: Attempt any \textbf{three} questions from this section.}\pts{3 $\times$ 15 = 45}

\begin{enumerate}[label=\textbf{\arabic*.},leftmargin=*,itemsep=10pt]

  \item Define Health. Explain various dimensions of health and factors affecting health.\pts{15}

  \item Define Physical Fitness. Discuss in detail the components of health-related and skill-related physical fitness.\pts{15}

  \item What is Exercise Physiology? Discuss the physiological response and adaptations to long-term aerobic exercise.\pts{15}

  \item Explain the principles of exercise prescription and describe how to design a balanced fitness program for youth.\pts{15}

\end{enumerate}

\bigskip

\section*{SECTION -- B}
\noindent \textit{Note: Attempt any \textbf{three} questions from this section.}\pts{3 $\times$ 5 = 15}

\begin{enumerate}[label=\textbf{\arabic*.},leftmargin=*,start=5,itemsep=8pt]

  \item Write short notes on Body Mass Index (BMI) and body composition.\pts{5}

  \item Differentiate between aerobic and anaerobic exercises with suitable examples.\pts{5}

  \item Explain the role of physical activity in the prevention of lifestyle diseases.\pts{5}

  \item Discuss warm-up and cool-down techniques in exercise management.\pts{5}

\end{enumerate}

\bigskip

\section*{SECTION -- C}
\noindent \textit{Note: Objective / Short Answer questions.}\pts{5 $\times$ 2 = 10}

\begin{enumerate}[label=\textbf{\arabic*.},leftmargin=*,start=9,itemsep=6pt]

  \item Define Wellness.\pts{2}
  \item What is Target Heart Rate (THR)?\pts{2}
  \item Name any two health-related physical fitness components.\pts{2}
  \item Define Obesity.\pts{2}
  \item Write the full form of WHO.\pts{2}

\end{enumerate}

\end{document}
